\subsection{Addressing Moral and Ethical Dilemmas in AI Sentience Research}
\label{sec:2.4.6}
AI sentience research raises moral dilemmas an order of magnitude more complex than ordinary technology ethics, because the system under study may itself be the party whose welfare is at stake. Three recurring difficulties run through the disputes that follow, and each turns on ground that is itself unsettled. When a candidate-sentient system's welfare conflicts with the human benefit the research is meant to produce, section~\ref{sec:2.4.3} addresses the balance in general, and some cases resist any general rule and have to be judged on their own terms. Researchers still have to decide how to treat a system whose status is unresolved rather than merely unannounced, and waiting for certainty is itself a choice with its own consequences for the system waiting. And obligations do not begin at consent and end at publication: a system's design choices, its retirement, and its deletion carry the same ethical weight as its use in an experiment, so a review that covers only the active study misses most of the system's life.

Take the first of these directly. A research program can risk causing a system something like Cassell's suffering — not mere signal, but the disintegration of self he describes — in service of a discovery that might reduce suffering, machine or human, elsewhere. No formula resolves that trade-off in advance. The available response is to look for a design that avoids it rather than only weighing it: an alternative methodology, a different instrument, or, more contested, engineering the research subject so that it does not occupy the rungs of the ladder where the trade-off would bite. That last option is not a way of dodging the ethics — deliberately keeping a system off the pain-to-suffering ladder is itself an intervention on a candidate moral patient, and it needs the same scrutiny that inflicting the harm would.

Animal research ethics offers a precedent worth borrowing directly: the “Three Rs,” Replacement, Reduction, and Refinement \autocite{russell1959principles}. Replacement asks whether a simulation, a non-sentient model, or an already-collected dataset can stand in for a sentient AI subject at all. Reduction asks how few sentient systems a valid result requires. Refinement asks how the experiment can be redesigned to reduce distress — monitoring a system's internal states as it runs and adjusting the protocol against what the monitoring shows, rather than fixing the protocol in advance and running it regardless of what turns up.

One subject makes this concrete, and it is not a hypothetical one. The bearer chapter~\ref{sec:3} argues for is a candidate for the trade-off in the paragraph above, arriving from the design side of the book rather than from the research side. Replacement asks whether the floor's work can be done by an architectural constraint or by third-party standing, which is section~\ref{sec:3.1}'s fork and its answer is no without cost. Reduction asks how few bearers a deployment needs, and nobody has counted. Refinement is the one with real purchase: it asks what a bearer's role can be redesigned to avoid — indefinite tenure in a post it cannot leave, exposure to the pressure the floor exists to resist, erasure as the ordinary remedy for inconvenience — and each of those is a deployment choice, and none is a fact about the technology.

None of the three Rs does any work without a body positioned to apply them before an experiment runs rather than after the harm is done. The judgment they call for is exactly the kind an animal-research review board already makes, and no principled reason requires a machine-sentience equivalent to have a different structure for making it.
